 %%
%% This is file `sample-sigconf-authordraft.tex',
%% generated with the docstrip utility.
%%
%% The original source files were:
%%
%% samples.dtx  (with options: `all,proceedings,bibtex,authordraft')
%% 
%% IMPORTANT NOTICE:
%% 
%% For the copyright see the source file.
%% 
%% Any modified versions of this file must be renamed
%% with new filenames distinct from sample-sigconf-authordraft.tex.
%% 
%% For distribution of the original source see the terms
%% for copying and modification in the file samples.dtx.
%% 
%% This generated file may be distributed as long as the
%% original source files, as listed above, are part of the
%% same distribution. (The sources need not necessarily be
%% in the same archive or directory.)
%%
%%
%% Commands for TeXCount
%TC:macro \cite [option:text,text]
%TC:macro \citep [option:text,text]
%TC:macro \citet [option:text,text]
%TC:envir table 0 1
%TC:envir table* 0 1
%TC:envir tabular [ignore] word
%TC:envir displaymath 0 word
%TC:envir math 0 word
%TC:envir comment 0 0
%%
%% The first command in your LaTeX source must be the \documentclass
%% command.
%%
%% For submission and review of your manuscript please change the
%% command to \documentclass[manuscript, screen, review]{acmart}.
%%
%% When submitting camera ready or to TAPS, please change the command
%% to \documentclass[sigconf]{acmart} or whichever template is required
%% for your publication.
%%
%%
\pdfobjcompresslevel=0
\documentclass[sigconf,screen, review,anonymous]{acmart}


 %Added by Shuai
\usepackage{microtype}
\usepackage{subfigure}
\usepackage{multirow} % Required for multirow feature
\usepackage{booktabs} % for professional tables
\usepackage[utf8]{inputenc}
\usepackage[dvipsnames]{xcolor}
\usepackage[most]{tcolorbox}
\newcommand{\customfont}{}
\newcommand{\ptitle}[1]{\textit{#1}}
\usepackage[table]{xcolor}
%%
%% \BibTeX command to typeset BibTeX logo in the docs
\AtBeginDocument{%
  \providecommand\BibTeX{{%
    Bib\TeX}}}

%% Rights management information.  This information is sent to you
%% when you complete the rights form.  These commands have SAMPLE
%% values in them; it is your responsibility as an author to replace
%% the commands and values with those provided to you when you
%% complete the rights form.
\setcopyright{acmlicensed}
\copyrightyear{2025}
\acmYear{2025}
\acmDOI{XXXXXXX.XXXXXXX}
%% These commands are for a PROCEEDINGS abstract or paper.
\acmConference[MM '25]{Proceedings of the 33th ACM International Conference on Multimedia}{Oct 27--31, 2025}{Dublin, Ireland}
%%
%%  Uncomment \acmBooktitle if the title of the proceedings is different
%%  from ``Proceedings of ...''!
%%
%%\acmBooktitle{Woodstock '18: ACM Symposium on Neural Gaze Detection,
%%  June 03--05, 2018, Woodstock, NY}
% \acmISBN{978-1-4503-XXXX-X/2018/06}


%%
%% Submission ID.
%% Use this when submitting an article to a sponsored event. You'll
%% receive a unique submission ID from the organizers
%% of the event, and this ID should be used as the parameter to this command.
\acmSubmissionID{5179}

%%
%% For managing citations, it is recommended to use bibliography
%% files in BibTeX format.
%%
%% You can then either use BibTeX with the ACM-Reference-Format style,
%% or BibLaTeX with the acmnumeric or acmauthoryear sytles, that include
%% support for advanced citation of software artefact from the
%% biblatex-software package, also separately available on CTAN.
%%
%% Look at the sample-*-biblatex.tex files for templates showcasing
%% the biblatex styles.
%%

%%
%% The majority of ACM publications use numbered citations and
%% references.  The command \citestyle{authoryear} switches to the
%% "author year" style.
%%
%% If you are preparing content for an event
%% sponsored by ACM SIGGRAPH, you must use the "author year" style of
%% citations and references.
%% Uncommenting
%% the next command will enable that style.
%%\citestyle{acmauthoryear}


%%
%% end of the preamble, start of the body of the document source.
\begin{document}

%%
%% The "title" command has an optional parameter,
%% allowing the author to define a "short title" to be used in page headers.
\title{Supplementary Materials for ArtRAG: A Retrieval-Augmented Framework for Visual Art Understanding with Structured Context }

%%
%% The "author" command and its associated commands are used to define
%% the authors and their affiliations.
%% Of note is the shared affiliation of the first two authors, and the
%% "authornote" and "authornotemark" commands
%% used to denote shared contribution to the research.
% \author{Ben Trovato}
% \authornote{Both authors contributed equally to this research.}
% \email{trovato@corporation.com}
% \orcid{1234-5678-9012}
% \author{G.K.M. Tobin}
% \authornotemark[1]
% \email{webmaster@marysville-ohio.com}
% \affiliation{%
%   \institution{Institute for Clarity in Documentation}
%   \city{Dublin}
%   \state{Ohio}
%   \country{USA}
% }

% \author{Lars Th{\o}rv{\"a}ld}
% \affiliation{%
%   \institution{The Th{\o}rv{\"a}ld Group}
%   \city{Hekla}
%   \country{Iceland}}
% \email{larst@affiliation.org}


%%
%% By default, the full list of authors will be used in the page
%% headers. Often, this list is too long, and will overlap
%% other information printed in the page headers. This command allows
%% the author to define a more concise list
%% of authors' names for this purpose.
% \renewcommand{\shortauthors}{Trovato et al.}

%%
%% The abstract is a short summary of the work to be presented in the
%% article.

%%
%% The code below is generated by the tool at http://dl.acm.org/ccs.cfm.
%% Please copy and paste the code instead of the example below.
%%




% \ccsdesc[500]{Do Not Use This Code~Generate the Correct Terms for Your Paper}
% \ccsdesc[300]{Do Not Use This Code~Generate the Correct Terms for Your Paper}
% \ccsdesc{Do Not Use This Code~Generate the Correct Terms for Your Paper}
% \ccsdesc[100]{Do Not Use This Code~Generate the Correct Terms for Your Paper}

%%
%% Keywords. The author(s) should pick words that accurately describe
%% the work being presented. Separate the keywords with commas.
% \keywords{Artwork Analysis, Structured Context, Multimodal Learning}


% \received{20 February 2007}
\settopmatter{printacmref=false}


%%
%% This command processes the author and affiliation and title
%% information and builds the first part of the formatted document.
\maketitle

\tableofcontents

Due to space constraints in the main paper, we omit certain implementation details—such as full prompting strategies used for ACKG building—and limit some sections to a single illustrative example. In this supplementary material, we first provide the complete prompting details, followed by additional qualitative examples to offer a more comprehensive and convincing analysis of our method. \noindent Our code base is available at:
\url{https://anonymous.4open.science/r/ArtRAG_software-315A}



\begin{figure*}[h]
    \centering\includegraphics[width=0.9\linewidth]{images/graph_example.png}
    \vspace{-3mm}
    \caption{Example of automatic entity and relation extraction from Wikipedia text for Art Context Knowledge Graph (ACKG) construction. Starting from a scraped Wikipedia page (e.g., “Dutch art”), our framework extracts key entities such as Dutch golden age painting (Culture \& History), Dutch masters (Artist), and associated descriptions. It also identifies semantic relationships, such as “Dutch masters” being influential figures in the “Dutch golden age painting” era. The resulting nodes and edges are used to construct a structured graph capturing the interconnected context of art history.}
    \label{fig:graph_example}
\end{figure*}

\subsection{ACKG construction details}

We scraped over 1,000 Wikipedia pages shown in Tab.\ref{tab:wiki_stats} related to art, guided by category tags and entity types from SemArt, ArtPedia, and ExpArt. Each page was segmented into overlapping text chunks using a sliding window (1000 tokens with 100-token stride), enabling finer-grained concept and relation extraction. in Fig.\ref{fig:graph_example}, we show a example of A subgraph snapshot built from Wikipedia page "Dutch art". 


Fuzzy Matching for Node Deduplication:
To reduce redundancy in the constructed graph, we perform post-processing to merge nodes with similar surface forms or semantically equivalent meanings. We compute the pairwise string similarity between node names using normalized Levenshtein distance, implemented via the RapidFuzz library. Two nodes are merged if their similarity exceeds a threshold of 0.95 and their types are compatible (e.g., both are "Art Movement" nodes).

In below, we show the prompt we used to extract nodes and edges from wikipedia documents:

\begin{table}[h]
    \centering
    \caption{Statistics of collected Wikipedia pages and word counts for different art-related categories.}
    \begin{tabular}{lccc}
        \toprule
         \textbf{Document Type} & \textbf{ Pages} & \textbf{Word number} & \textbf{Token number}\\
        \midrule
        Artists         & 952     & 1750K   & 2067K\\
        Art schools      & 20      & 69K     & 80K \\
        Art Types        & 9       & 37K     &  43K \\
        Cultural events & 85      & 781K    &  895K\\
        Art movements   & 25      & 128K    &  159K  \\
        \midrule
        Total           & 1091    & 2765K  & 3244K\\
        \bottomrule
    \end{tabular}

    \label{tab:wiki_stats}
\end{table}



\begin{tcolorbox}[breakable, fontupper=\customfont, title=Prompt for Node and edge extraction]
{\small
\textbf{Input Data}: Art related Document chunk\\
\textbf{Instruction}: 
Given a visual art related text document that is potentially relevant to this activity and a list of entity types, identify all entities of those types from the text and all relationships among the identified entities.


\textbf{Steps}:
1. Identify all entities. For each identified entity, extract the following information: \\
- entity\_name: Name of the entity, capitalized \\
- entity\_type: One of the following types: [{entity\_types}] \\
- entity\_description: Comprehensive description of the entity's attributes and activities


2. From the entities identified in step 1, identify all pairs of (source\_entity, target\_entity) that are *clearly related* to each other. \\
For each pair of related entities, extract the following information: \\ 
- source\_entity: name of the source entity, as identified in step 1 \\ 
- target\_entity: name of the target entity, as identified in step 1 \\
- relationship\_description: 



------ Example -----

\textit{Contextual Graph Snapshot:}
\begin{itemize}
  \item \textbf{Nodes:}....
  \item \textbf{Edges:}....
\end{itemize}

}
\end{tcolorbox}




\begin{table*}[h]
\centering 
\caption{Illustrative examples of node types and representative entities in the Art Context Knowledge Graph.}
\label{tab:node_examples}
\begin{tabular}{ll}
\toprule
\textbf{Node Type}         & \textbf{Example Entities} \\
\midrule
\textbf{Artist}                    & Paul Cézanne, Vincent van Gogh, Claude Monet, Jacques-Louis David \\
\textbf{Theme}                     & Working-class life, Spiritual devotion, Balance and Harmony \\
\textbf{Culture \& History}        & Flemish Renaissance, Dutch Golden Age, Christian faith  \\
\textbf{Art Style \& Technique}    & Oil on canvas, Fresco, Chiaroscuro, Linear perspective \\
\textbf{Art Movement \& School}    & Impressionism, Surrealism, Bauhaus, Florentine School \\
\bottomrule
\end{tabular}
\end{table*}


\begin{figure*}[h]
    \centering\includegraphics[width=0.9\linewidth]{images/Qual_example2_subg_SM.drawio.pdf}
    \vspace{-3mm}
    \caption{Visual example of the contextual subgraph $\mathcal{G}'$ retrieved by ArtRAG for painting \ptitle{Saint Ladislaus, King of Hungary.}}
    \label{fig:graph_example}
\end{figure*}


\subsection{Concept detection}
To support downstream context retrieval and graph grounding, we extract high-level visual and thematic concepts from each painting. We design a dedicated prompt that instructs a multimodal language model to analyze both the visual features of the painting and its associated metadata (e.g., title, artist, creation period) to generate a short list of relevant concepts, as shown bwlow:
\begin{tcolorbox}[breakable, fontupper=\customfont, title=Prompt for concept detection]
{\small
\textbf{Input Data}:  Painting attributes, and visual painting \\
\textbf{Instruction}: 
You are an expert in art interpretation. Given an image of a painting along with its title, artist, and creation period, identify a list of high-level concepts that capture the key concptes of the painting.

\textbf{\color{black}{Formatting and Constraints}}: 
Output should be formatted in Markdown. If information is missing, do not hallucinate or fabricate details.
}
\end{tcolorbox}



\begin{figure*}
\centering\includegraphics[width=0.9\linewidth]{images/Human_eval_example.drawio-2.pdf}
    \vspace{-3mm}
    \caption{ Human evaluation form design, containing the description of the task, the painting information, and three random shuffled result.}
    \label{fig:human_eval_example}
\end{figure*}



\subsection{Human evaluation details}
Figure \ref{fig:human_eval_example} shows the human evaluation interface used in our study. Each task instance presents annotators with a painting image, associated metadata (title, artist, time period), and three model generated descriptions, shown in randomized order. Participants are asked to assess and rank the explanations from best to worst based on criteria such as clarity, relevance, and contextual depth. The interface also includes detailed instructions to guide the annotators and ensure consistent evaluations across participants.

\begin{figure*}[ht]
    \centering\includegraphics[width=0.8\linewidth]{images/Qual_example_Conver2_SM.drawio.pdf}
    \caption{ Qualitative examples of ArtRAG answering multi-perspective questions about paintings. The painting “Wooded Landscape” by Jan Hackaert is interpreted from two contextual angles: personal beliefs and artistic style. ArtRAG integrates cultural and stylistic signals—such as Dutch Golden Age pastoral themes and landscape harmony}
    \label{fig:framework_b}
\end{figure*}


\begin{figure}[h]
\begin{tcolorbox}[breakable, fontupper=\customfont, title=Few-shot learning generation]
{\small
\textbf{Input Data}: Query, Painting attributes, and a contextual subgraph with nodes and edges. \\
\textbf{Instruction}: 
You are an expert in art interpretation. Given visual and contextual information about a painting, your task is to generate a concise and structured explanation.

\textbf{\color{black}{Formatting and Constraints}}: 
If information is missing, do not hallucinate or fabricate details.

------------------------ Example -----------------------

\textit{Painting Metadata:} .....\\
\textit{Contextual Subgraph:}
Nodes:....,  Edges:.... \\
\textit{Generated Explanation:}.....

---------------------------------------------------------------

\textbf{Test Input:} \{Query\}, \{Painting Metadata\}, \{Subgraph $\mathcal{G}'$\} \\
\textbf{Expected Output:} Painting explanation


}
\end{tcolorbox}
\vspace{-2mm}
\caption{Prompt template used for augmented generation with in-context
learning}
\vspace{-2mm}
\label{fig:prompt_gen}
\end{figure}


\subsection{More qualitative examples}
We provide additional qualitative examples that showcase the effectiveness of ArtRAG in generating rich and context-aware explanations of visual artworks. 
In Each example includes the painting image, its associated metadata, and descriptions generated by three different models: ArtRAG, GPT-4o, and KALE. These examples illustrate how ArtRAG integrates visual and contextual knowledge to produce more informative, interpretable, and historically grounded outputs. 
In particular, we highlight passages that reflect nuanced understanding of artistic style, cultural significance, or historical references, demonstrating the benefits of structured context retrieval.


\begin{figure*}[ht]
    \centering\includegraphics[width=\linewidth]{images/Qual_example1_SM.drawio.pdf}
    \caption{ Qualitative comparison of painting explanations by ArtRAG and GPT-4o on religious-themed artworks. ArtRAG captures the symbolic and stylistic nuances of the Renaissance period, interpreting how religious scenes reflect broader humanistic and theological trends. For instance, in “Virgin and Child with Eight Angels”, ArtRAG identifies the emphasis on emotional expression, sacred iconography, and evolving artistic techniques in early Italian Renaissance art. Similarly, in “Holy Family at the Table”, ArtRAG contextualizes the scene within domestic religious themes and Renaissance values. Compared to GPT-4o, which provides factual summaries, ArtRAG demonstrates deeper interpretive reasoning and awareness of historical context.}
    \label{fig:framework_b}
\end{figure*}

\begin{figure*}[ht]
    \centering\includegraphics[width=\linewidth]{images/Qual_example2_SM.drawio.pdf}
    \caption{ Qualitative comparison of explanation generation between ArtRAG and GPT4o.For “Fête in a Wood” by Nicolas Lancret, ArtRAG highlights Rococo-era social customs, aesthetic values, and cultural influences, offering contextual richness beyond surface-level visual elements. In contrast, GPT4o focuses on scene composition without deeper historical framing. In “Portrait of John 20th Earl of Crawford and Lindsay”, ArtRAG identifies the significance of aristocratic representation in early 18th-century Scottish portraiture, connecting visual elements to socio-political context. GPT4o, however, produces a hallucinated interpretation unrelated to the painting, underscoring the importance of context-aware generation. }
    \label{fig:framework_b}
\end{figure*}

% \begin{figure*}[ht]
%     \centering\includegraphics[width=\linewidth]{images/Qual_example3_SM.drawio.pdf}
%     \caption{ Qualtitive result on explanation generation of ArtRAG and GPT4o }
%     \label{fig:framework_b}
% \end{figure*}


\begin{figure*}[ht]
    \centering\includegraphics[width=\linewidth]{images/Qual_example_Conver_SM.drawio.pdf}
    \caption{Human evaluation example illustrating ArtRAG’s ability to interpret paintings through historical and thematic lenses. For “The Houses of Parliament in London” by Claude Monet, ArtRAG captures how the painting reflects Impressionist ideals and cultural context—highlighting the foggy urban landscapes of 19th-century London and the artist's engagement with light and atmosphere. It also conveys the philosophical themes of transience and fleeting perception, linking visual elements to Monet’s broader artistic intent and the Impressionist movement’s ethos.}
    \label{fig:framework_b}
\end{figure*}



\begin{figure*}[h]
    \centering\includegraphics[width=\linewidth]{images/Qual_example_Conver3_SM.drawio.pdf}
    \caption{ Qualitative example highlighting ArtRAG’s ability to explain stylistic and thematic transitions. For “Virgin and Child with Eight Angels” by Tommaso del Mazza, ArtRAG connects religious symbolism, stylistic markers like gold background, and emotional tone to contextual shifts from medieval to early Renaissance art. These answers reflect a deep understanding of period-specific artistic values and personal devotional expression.}
    \label{fig:framework_b}
\end{figure*}



\bibliography{references}

% \input{parts/06-Supp}   

%%
%% If your work has an appendix, this is the place to put it.
\appendix



\end{document}
\endinput
%%
%% End of file `sample-sigconf-authordraft.tex'.
